\chapter{Chapitre IV : intégration et dérivation numériques}

% ------------------------------------------------------------------
\section{Exercice IV.1 : dérivation à partir de l'interpolation quadratique}

\Enonce{Établir la formule de dérivation numérique à partir de l'interpolation
quadratique.}

\Pourquoi On ne connaît $f$ qu'en quelques points. On la remplace par le
polynôme de degré 2 qui passe par trois points, et on dérive ce polynôme. Avec
trois points au lieu de deux, l'erreur passe de l'ordre $h$ (formule classique)
à l'ordre $h^2$.

\Etapes
\Etape{1} Points équidistants $x_1 = x - h$, $x_2 = x$, $x_3 = x + h$. On part de
$f(t) \approx f(x_1)f_1(t) + f(x_2)f_2(t) + f(x_3)f_3(t)$ et on dérive les
trois polynômes de base. Par exemple
$f_1(t) = (t - x_2)(t - x_3)/\bigl((x_1 - x_2)(x_1 - x_3)\bigr)$ donne
$f_1'(t) = (2t - x_2 - x_3)/\bigl((x_1 - x_2)(x_1 - x_3)\bigr)$.

\Etape{2} Au point $t = x_2$ :
\[
\begin{aligned}
  f_1'(x_2) &= \frac{x_2 - x_3}{(x_1 - x_2)(x_1 - x_3)} = \frac{-h}{(-h)(-2h)} = -\frac{1}{2h},\\
  f_2'(x_2) &= \frac{2x_2 - x_1 - x_3}{(x_2 - x_1)(x_2 - x_3)} = 0,\\
  f_3'(x_2) &= \frac{x_2 - x_1}{(x_3 - x_1)(x_3 - x_2)} = \frac{h}{(2h)(h)} = \frac{1}{2h}.
\end{aligned}
\]
D'où la formule centrée
\[
  f'(x) \approx \frac{f(x+h) - f(x-h)}{2h} .
\]

\Etape{3} Erreur. Les développements de Taylor à l'ordre 3,
$f(x \pm h) = f(x) \pm hf'(x) + \tfrac{h^2}{2}f''(x) \pm \tfrac{h^3}{6}f'''(\xi_\pm)$,
donnent par différence
\[
  \frac{f(x+h) - f(x-h)}{2h} = f'(x) + \frac{h^2}{12}\bigl(f'''(\xi_+) + f'''(\xi_-)\bigr)
  = f'(x) + \frac{h^2}{6}f'''(\xi),
\]
la dernière égalité venant du théorème des valeurs intermédiaires ($f'''$
continue). On retrouve l'erreur du cours $E = -\tfrac{1}{6}h^2f'''(\xi)$ :
$f'(x) = \frac{f(x+h) - f(x-h)}{2h} - \frac{h^2}{6}f'''(\xi)$.

\Etape{4} Au point $t = x_1$ (dérivée en bout d'intervalle, utile au bord d'un
tableau), le même calcul donne $f_1'(x_1) = -3/(2h)$, $f_2'(x_1) = 2/h$,
$f_3'(x_1) = -1/(2h)$, c'est-à-dire la seconde formule du cours
$f'(x) \approx \bigl[-3f(x) + 4f(x+h) - f(x+2h)\bigr]/(2h)$, d'erreur
$+\tfrac13h^2f'''(\xi)$.

\Refs \BF{§~4.1, formules à trois points (4.4) et (4.5), p.~174--175}.

% ------------------------------------------------------------------
\section{Exercice IV.2 : vitesses à partir d'un tableau de positions}

\Enonce{$t = 1{,}0$ ; $1{,}1$ ; \dots ; $1{,}5$ et
$x = 1{,}6487$ ; $1{,}7333$ ; $1{,}8221$ ; $1{,}9155$ ; $2{,}0138$ ; $2{,}1170$.
Calculer les vitesses à $t = 1{,}05$ ; $1{,}15$ ; $1{,}25$ ; $1{,}35$ et $1{,}45$.}

\Pourquoi Les instants demandés sont les \emph{milieux} des intervalles du
tableau. La formule centrée de l'exercice IV.1, appliquée avec le pas $h/2$
autour du milieu $t + h/2$, n'utilise justement que $x(t)$ et $x(t+h)$ :
\[
  v\Bigl(t + \frac h2\Bigr) \approx \frac{x(t+h) - x(t)}{h},
  \qquad\text{erreur } -\frac{(h/2)^2}{6}x'''(\xi) = -\frac{h^2}{24}x'''(\xi).
\]
C'est la deuxième formule classique du cours ($f(x + h/2) - f(x - h/2)$ divisé
par $h$). On obtient donc une précision d'ordre $h^2$ avec les seules données.

\Etapes
\Etape{1} Différences successives des positions :
$0{,}0846$ ; $0{,}0888$ ; $0{,}0934$ ; $0{,}0983$ ; $0{,}1032$.

\Etape{2} Division par $h = 0{,}1$ :
\[
  v(1{,}05) \approx 0{,}846, \quad v(1{,}15) \approx 0{,}888, \quad
  v(1{,}25) \approx 0{,}934, \quad v(1{,}35) \approx 0{,}983, \quad
  v(1{,}45) \approx 1{,}032 .
\]

\Etape{3} Précision. Les positions sont données à $5\times 10^{-5}$ près ; une
différence divisée par $h = 0{,}1$ peut donc être fausse d'environ
$2 \times 5\times 10^{-5}/0{,}1 = 10^{-3}$. Les vitesses ne sont fiables qu'à
$10^{-3}$ près environ : c'est l'erreur d'arrondi, qui croît quand $h$ diminue,
alors que l'erreur de troncature décroît \BF{p.~179}.

\Verif On constate que les données coïncident avec $e^{t/2}$ à
$4{,}7\times 10^{-5}$ près (calcul fait sur les six points). Si c'est bien la
loi du mouvement, la vitesse exacte est $\tfrac12e^{t/2}$, soit
$0{,}8452$ ; $0{,}8886$ ; $0{,}9341$ ; $0{,}9820$ ; $1{,}0324$ : les écarts
avec nos valeurs restent inférieurs à $10^{-3}$, comme prévu à l'étape 3.

\Refs \BF{§~4.1, formule (4.5) p.~175 ; erreurs d'arrondi p.~178--179}.

% ------------------------------------------------------------------
\section{\texorpdfstring{Exercice IV.3 : trapèzes à la main pour $\int_1^2 \dd x/x$}{Exercice IV.3 : trapèzes à la main pour l'intégrale de 1/x}}

\Enonce{Calculer à la main $I = \int_1^2 \dd x/x$ pour $h = 0{,}2$, puis pour
$h = 0{,}1$. Comparer à la valeur exacte.}

\Pourquoi La valeur exacte est connue, $I = \ln 2 = 0{,}693147$ : l'exercice
permet de vérifier la formule composite des trapèzes et la loi d'erreur en
$h^2$ établie dans le cours, $E = -\tfrac{h^2}{12}(b-a)f''(\eta)$.

\Etapes
\Etape{1} $h = 0{,}2$ ($n = 5$). Valeurs de $f(x) = 1/x$ :
$1$ ; $0{,}833333$ ; $0{,}714286$ ; $0{,}625$ ; $0{,}555556$ ; $0{,}5$. Somme
des valeurs intérieures : $2{,}728175$.
\[
  T_{0,2} = \frac{0{,}2}{2}\bigl[1 + 2 \times 2{,}728175 + 0{,}5\bigr]
          = 0{,}1 \times 6{,}956349 = 0{,}695635 .
\]

\Etape{2} $h = 0{,}1$ ($n = 10$). Valeurs intérieures : $0{,}909091$ ;
$0{,}833333$ ; $0{,}769231$ ; $0{,}714286$ ; $0{,}666667$ ; $0{,}625$ ;
$0{,}588235$ ; $0{,}555556$ ; $0{,}526316$, de somme $6{,}187714$.
\[
  T_{0,1} = 0{,}05\bigl[1 + 2 \times 6{,}187714 + 0{,}5\bigr] = 0{,}693771 .
\]

\Etape{3} Comparaison. Erreurs : $T_{0,2} - \ln 2 = 0{,}002488$ et
$T_{0,1} - \ln 2 = 0{,}000624$. Leur rapport vaut $3{,}99 \approx 4$ : diviser
$h$ par 2 divise l'erreur par 4, c'est la loi en $h^2$. Les erreurs sont
positives, ce qu'annonce la formule car $f''(x) = 2/x^3 > 0$. Majoration : sur
$[1,2]$, $\abs{f''} \le 2$, d'où $\abs{E} \le h^2/6$, soit $0{,}0067$ et
$0{,}0017$ : elle est bien respectée.

\Etape{4} Prolongement (Richardson--Romberg). Comme l'erreur est en $h^2$, la
combinaison $(4T_{0,1} - T_{0,2})/3 = 0{,}693150$ élimine le terme principal :
l'erreur tombe à $3\times 10^{-6}$.

\Refs \BF{th.~4.5, p.~205 (trapèzes composites) ; intégration de Romberg,
§~4.5, p.~211--216}.

% ------------------------------------------------------------------
\section{Exercice IV.4 : trapèzes sans tableaux}

\Enonce{Écrire un autre algorithme qui n'utilise pas les tableaux.}

\Pourquoi Dans l'algorithme du cours, chaque valeur $f_k$ n'est utilisée
qu'une fois : il est inutile de la stocker. On l'ajoute à la somme dès qu'elle
est calculée. La mémoire utilisée ne dépend plus de $n$. On calcule les nœuds
par $x = a + kh$ plutôt que par additions répétées $x \leftarrow x + h$, pour
ne pas accumuler les erreurs d'arrondi.

\Etapes
\Etape{1} Initialiser la somme avec les deux extrémités (poids $1/2$).

\Etape{2} Ajouter les $n-1$ valeurs intérieures (poids $1$).

\Etape{3} Multiplier par $h$.

\begin{algo}[ : trapèzes composites sans tableaux]
  \Require la fonction $f$, $a$, $b$, $n$
  \State $h \leftarrow (b-a)/n$
  \State $S \leftarrow \bigl(f(a) + f(b)\bigr)/2$
  \For{$k$ allant de $1$ à $n-1$}
    \State $S \leftarrow S + f(a + kh)$
  \EndFor
  \State $S \leftarrow h\,S$
  \Print $S$
\end{algo}

\Verif Cet algorithme est programmé dans \fichier{ex11_trapeze.pas} et
\fichier{ex11_trapeze.f90} (exercice 11 de la fin du document) ; appliqué à
$1/x$ sur $[1,2]$, il redonne les valeurs de l'exercice IV.3.

\Refs \BF{p.~205} (même organisation dans l'algorithme 4.1).

% ------------------------------------------------------------------
\section{Exercice IV.5 : algorithme de Simpson}

\Enonce{Écrire un algorithme pour la méthode de Simpson.}

\Pourquoi La formule de Simpson s'applique sur des paires d'intervalles (un
polynôme de degré 2 sur trois points) : le nombre $n$ de sous-intervalles doit
donc être \emph{pair}. Les poids sont $1, 4, 2, 4, \dots, 2, 4, 1$ : $4$ aux
nœuds d'indice impair (milieux des paires), $2$ aux nœuds pairs intérieurs
(extrémités communes à deux paires). L'erreur composite vaut
$-\tfrac{b-a}{180}h^4f^{(4)}(\mu)$ \BF{th.~4.4, p.~204} : la méthode est
exacte pour les polynômes de degré 3.

\Etapes
\Etape{1} Vérifier que $n$ est pair.

\Etape{2} Accumuler séparément la somme $S_1$ des valeurs aux nœuds impairs et
la somme $S_2$ des valeurs aux nœuds pairs intérieurs (c'est ce que fait
l'algorithme 4.1 de Burden et Faires avec ses variables $XI_1$ et $XI_2$).

\Etape{3} $I \approx \tfrac h3\bigl[f(a) + 4S_1 + 2S_2 + f(b)\bigr]$.

\begin{algo}[ : Simpson composite]
  \Require $f$, $a$, $b$, $n$ (pair)
  \If{$n$ impair} \Print \og $n$ doit être pair \fg{} \Stop \EndIf
  \State $h \leftarrow (b-a)/n$ ; $S_1 \leftarrow 0$ ; $S_2 \leftarrow 0$
  \For{$k$ allant de $1$ à $n-1$}
    \State $x \leftarrow a + kh$
    \If{$k$ impair} \State $S_1 \leftarrow S_1 + f(x)$
    \Else \State $S_2 \leftarrow S_2 + f(x)$
    \EndIf
  \EndFor
  \State $I \leftarrow \frac{h}{3}\bigl(f(a) + 4S_1 + 2S_2 + f(b)\bigr)$
  \Print $I$
\end{algo}

\Verif Voir l'exercice IV.7 (calcul de $\pi$) et le programme
\fichier{ex12_mc.pas} (exercice 12 de la fin du document).

\Refs \BF{alg.~4.1, p.~205 ; th.~4.4, p.~204}.

% ------------------------------------------------------------------
\section{\texorpdfstring{Exercice IV.6 : formule d'intégration de degré 3 sur $a$, $c$, $d$, $b$}{Exercice IV.6 : formule d'intégration de degré 3 sur a, c, d, b}}

\Enonce{Établir une formule d'intégration numérique qui utilise une
interpolation de degré 3 et qui interpole les points $a$, $b$,
$c = (a+b)/4$ et $d = 3(a+b)/4$.}

\Attention Les points $(a+b)/4$ et $3(a+b)/4$ ne sont dans $[a,b]$, aux
quarts de l'intervalle, que si $a = 0$. On interprète donc l'énoncé comme
$c = a + (b-a)/4$ et $d = a + 3(b-a)/4$, ce qui coïncide avec lui lorsque
$a = 0$.

\Pourquoi Intégrer le polynôme d'interpolation de degré 3 revient à chercher
des poids $w_k$ tels que $\int_a^b f \approx \sum_k w_kf(x_k)$ soit exact pour
$1, x, x^2, x^3$. Les quatre nœuds étant symétriques par rapport au milieu de
$[a,b]$, les poids le sont aussi : il ne reste que deux inconnues.

\Etapes
\Etape{1} Changement de variable $x = a + (b-a)t$, $t \in [0,1]$ :
$\int_a^b f(x)\,\dd x = (b-a)\int_0^1 f\bigl(a + (b-a)t\bigr)\dd t$. Les nœuds
deviennent $t = 0$, $\tfrac14$, $\tfrac34$, $1$.

\Etape{2} Par symétrie, poids $\alpha$ en $0$ et $1$, et $\beta$ en $\tfrac14$ et
$\tfrac34$. Exactitude pour $1$ : $2\alpha + 2\beta = 1$. Exactitude pour
$t^2$ : $\alpha(0 + 1) + \beta\bigl(\tfrac1{16} + \tfrac9{16}\bigr) = \tfrac13$,
soit $\alpha + \tfrac58\beta = \tfrac13$.

\Etape{3} De la première, $\alpha = \tfrac12 - \beta$ ; en reportant :
$\tfrac12 - \tfrac38\beta = \tfrac13$, d'où $\beta = \tfrac49$ et
$\alpha = \tfrac1{18}$. L'exactitude pour $t$ et $t^3$ est automatique (les
fonctions $t - \tfrac12$ et $(t - \tfrac12)^3$ sont impaires par rapport au
milieu, et les poids sont symétriques).

\Etape{4} Formule obtenue :
\[
  \int_a^b f(x)\,\dd x \approx \frac{b-a}{18}\Bigl[f(a) + 8f(c) + 8f(d) + f(b)\Bigr] .
\]

\Etape{5} Erreur. Pour $t^4$, la formule donne
$\tfrac49\bigl(\tfrac1{256} + \tfrac{81}{256}\bigr) + \tfrac1{18} = \tfrac{19}{96}$
au lieu de $\tfrac15$ : l'écart (exact moins formule) vaut $\tfrac1{480}$. Comme
$(t^4)^{(4)} = 24$, l'erreur est, pour $f^{(4)}$ à peu près constante,
\[
  E \approx \frac{(b-a)^5}{11\,520}\,f^{(4)} .
\]
Le noyau de Peano de cette formule change légèrement de signe (calcul
numérique) : on ne peut pas écrire rigoureusement $E = K f^{(4)}(\xi)$, mais
on a la majoration $\abs{E} \le 9{,}13\times 10^{-5}(b-a)^5\max\abs{f^{(4)}}$.

\Verif Sur $\int_0^1 e^{x}\dd x = e - 1$ : l'erreur vaut $1{,}44\times 10^{-4}$,
proche de la prévision $e^{0{,}5}/11\,520 = 1{,}43\times 10^{-4}$. À titre de
comparaison, la formule de Newton-Cotes à quatre points équidistants
(\og 3/8 de Simpson \fg, erreur $-\tfrac{3}{80}h^5f^{(4)}$ avec
$h = (b-a)/3$, soit $(b-a)^5/6480$) donne $-2{,}58\times 10^{-4}$ : la
formule trouvée est environ deux fois plus précise.

\Refs \BF{§~4.3 : degré de précision, p.~196 ; formules de Newton-Cotes
fermées et règle des 3/8, th.~4.2, p.~196--197}.

% ------------------------------------------------------------------
\section{\texorpdfstring{Exercice IV.7 : Simpson à la main pour $\int_0^1 4/(1+x^2)\,\dd x$}{Exercice IV.7 : Simpson à la main pour l'intégrale de 4/(1+x²)}}

\Enonce{Utiliser la méthode de Simpson pour calculer manuellement
$I = \int_0^1 \frac{4}{1+x^2}\dd x$ avec $h = 0{,}1$.}

\Pourquoi Une primitive est $4\arctan x$, donc $I = 4\arctan 1 = \pi$. On peut
ainsi mesurer l'erreur exacte de la méthode.

\Etapes
\Etape{1} $n = 10$ (pair). Valeurs $f_k = 4/(1 + x_k^2)$ :
\begin{center}
\begin{tabular}{c|cccccc}
  $k$ & 0 & 1 & 2 & 3 & 4 & 5\\ \hline
  $f_k$ & 4 & 3,960396 & 3,846154 & 3,669725 & 3,448276 & 3,2\\[4pt]
  $k$ & 6 & 7 & 8 & 9 & 10 & \\ \hline
  $f_k$ & 2,941176 & 2,684564 & 2,439024 & 2,209945 & 2 &
\end{tabular}
\end{center}

\Etape{2} Sommes : indices impairs
$S_1 = 3{,}960396 + 3{,}669725 + 3{,}2 + 2{,}684564 + 2{,}209945 = 15{,}724629$ ;
indices pairs intérieurs
$S_2 = 3{,}846154 + 3{,}448276 + 2{,}941176 + 2{,}439024 = 12{,}674631$.

\Etape{3}
\[
  I \approx \frac{0{,}1}{3}\bigl[4 + 4 \times 15{,}724629 + 2 \times 12{,}674631 + 2\bigr]
    = \frac{0{,}1}{3} \times 94{,}247778 = 3{,}1415926 .
\]

\Etape{4} Comparaison : $\pi = 3{,}14159265$ ; l'erreur vaut
$-4\times 10^{-8}$. La majoration $\frac{b-a}{180}h^4\max\abs{f^{(4)}}$ avec
$\max_{[0,1]}\abs{f^{(4)}} = 96$ (atteint en $x = 0$) vaut $5{,}3\times
10^{-5}$ : elle est respectée, mais très pessimiste ici.

\Verif \fichier{calculs.py} : $S = 3{,}1415926139$, erreur
$-3{,}965\times 10^{-8}$.

\Refs \BF{th.~4.4, p.~204 ; alg.~4.1, p.~205}.
